\FifthArticle{17}
Das Verfahren, das zur Bestimmung der imaginären Wurzeln angewandt werden muss, unterscheidet sich ganz von dem in den vorangehenden Artikeln gelehrten. Im Allgemeinen ist die indirekte Bestimmung imaginärer Wurzeln weit schwieriger als die reeller Wurzeln, weil jene in einem unendlichen Gebiet von zwei Dimensionen gesucht werden müssen, diese nur in einem Unendlichen von einer Dimension. Gerade deshalb verdient ein sehr umfassender besonderer Fall eine eigene Ausführung, in dem man jene Schwierigkeit umgehen und die Frage in dasselbe Gebiet versetzen kann, zu dem die Suche nach reellen Wurzeln gehört. Einen solchen Fall bieten die Gleichungen mit drei Gliedern.

Da sich die Methode mit gleicher Leichtigkeit anwenden lässt, gleichgültig ob die Koeffizienten der Gleichung reell oder imaginär sind, lege ich sofort die allgemeine Form der Gleichung zugrunde:
\[
\mathrm{X}=x^{m+n}+e(\cos\varepsilon+i\sin\varepsilon)x^m+f(\cos\varphi+i\sin\varphi)=0
\]
Dabei bedeuten $e$ und $f$ positive Größen. Für einen reellen, positiven oder negativen Koeffizienten ist der betreffende Winkel ($\varepsilon$ oder $\varphi$) dann entweder $0$ oder $180^{0}$. Die Voraussetzung, dass $m$ und $n$ keinen gemeinsamen Teiler haben, kann auch hier ohne Beeinträchtigung der Allgemeinheit beibehalten werden. Eine imaginäre Wurzel $x=t+iu$, die der Gleichung genügt, setzt man in die Form $r(\cos\rho+i\sin\rho)$. Für unseren gegenwärtigen Zweck ist es
\SourceLeaf{104}
vorteilhafter, die sonst übliche Bedingung, dass $r$ positiv sein soll, hier nicht zu stellen, sondern stattdessen zu verlangen, dass $\rho$ immer zwischen den Grenzen $0$ und $180^{0}$ genommen wird. Sind beide Koeffizienten der Gleichung reell, so lässt sich der Umfang der Werte von $\rho$ noch weiter auf die Hälfte einschränken. Da die imaginären Wurzeln einer solchen Gleichung bekanntlich paarweise zusammengehören, etwa $t+iu$ und $t-iu$, fällt für die eine Wurzel jedes Paares der Wert von $\rho$ offenbar zwischen $0$ und $90^{0}$. Nur diese braucht man durch das indirekte Verfahren zu bestimmen; die andere folgt daraus von selbst, indem man $\rho$ mit $180^{0}-\rho$ und $r$ mit $-r$ vertauscht.

\FifthArticle{18}
Das Wesen der Methode besteht darin, eine Gleichung aufzustellen, die nur $\rho$, ohne $r$, enthält. Um sie zu erhalten, bringt man die Gleichung $\mathrm{X}=0$ durch Division durch ihr erstes Glied in die Form
\[
1+e(\cos\varepsilon+i\sin\varepsilon)x^{-n}+f(\cos\varphi+i\sin\varphi)x^{-m-n}=0
\]
oder
\[
\begin{aligned}&1+er^{-n}(\cos(n\rho-\varepsilon)-i\sin(n\rho-\varepsilon))\\
&\quad+fr^{-m-n}(\cos[(m+n)\rho-\varphi]-i\sin[(m+n)\rho-\varphi])=0\end{aligned}
\]
Da sich hier die imaginären Teile aufheben müssen, hat man\nobreak\hspace{.5em}(I)
\[
r^m=-\frac{f\sin((m+n)\rho-\varphi)}{e\sin(n\rho-\varepsilon)}
\]
Auf ähnliche Weise erhält man, wenn man die Gleichung $\mathrm{X}=0$ durch ihr zweites oder drittes Glied dividiert und beachtet, dass sich in beiden Fällen die imaginären Teile der neuen Gleichungen aufheben müssen, die Gleichungen
\[
\left.
\begin{array}{rcl}
r^{m+n}&=&\displaystyle\frac{f\sin(m\rho+\varepsilon-\varphi)}{\sin(n\rho-\varepsilon)}\\[1.2ex]
r^n&=&\displaystyle-\frac{e\sin(m\rho+\varepsilon-\varphi)}{\sin((m+n)\rho-\varphi)}
\end{array}
\right\}\quad\mathrm{(I)}
\]
Man sieht, dass sich jede der drei Gleichungen (I) auch schon aus der Verbindung der beiden anderen ableiten lässt. Eliminiert man aber $r$ aus der Verbindung zweier Gleichungen, so erhält man (II)
\SourceLeaf{105}
\[
\lambda=(-1)^{m+n}\frac{\sin(m\rho+\varepsilon-\varphi)^m\sin(n\rho-\varepsilon)^n}{\sin((m+n)\rho-\varphi)^{m+n}}
\]
wobei zur Abkürzung wie oben
\[
\frac{f^n}{e^{m+n}}=\lambda
\]
gesetzt ist. Aus dieser Gleichung sind die verschiedenen Werte von $\rho$ zu bestimmen. Den Wert von $r$, der jedem Wert von $\rho$ entspricht, findet man dann aus einer der Gleichungen (I), dem Betrag nach am besten aus der zweiten. Ist jedoch $m+n$ gerade, muss noch eine der beiden anderen Gleichungen zur Entscheidung des Vorzeichens herangezogen werden.

\FifthArticle{19}
Die Gleichung II wird man auf indirektem Weg immer mit Leichtigkeit lösen können. Dazu wird auch die Berücksichtigung der folgenden Bemerkungen beitragen.

1) Die Werte von $\rho$ liegen zwischen $0$ und $180^{0}$. Sind die Koeffizienten der vorgegebenen Gleichung reell, so braucht man nur die halbe Anzahl einzeln aufzusuchen, nämlich die zwischen $0$ und $90^{0}$ liegenden.

2) In beiden Fällen teilt man zunächst das betreffende Intervall in die verschiedenen Unterabschnitte, die durch die Vorzeichenwechsel der Werte der auf der rechten Seite der Gleichung II stehenden Funktion von $\rho$ entstehen. Die Übergangswerte von $\rho$ können offenbar nur dort liegen, wo einer der Winkel $m\rho+\varepsilon-\varphi,\ n\rho-\varepsilon,\ (m+n)\rho-\varphi$ durch $180^{0}$ teilbar ist und daher die Funktion selbst entweder $0$ oder unendlich wird. Diejenigen dieser Unterabschnitte, in denen der Wert der Funktion negativ wird, sind dann von selbst von der weiteren Untersuchung ausgeschlossen.

3) Falls man nicht schon auf anderen Wegen Näherungswerte von $\rho$ erhalten kann, wird man sich die indirekte Suche in den geeigneten Intervallen sehr erleichtern, indem man die ersten Versuche mit abgekürzten Tafeln mit wenigen Ziffern durchführt, ähnlich wie es die Beispiele von Artikel 16 zeigen. In manchen Fällen mag es bequem sein,
\SourceLeaf{106}
dafür zunächst nur die Logarithmen der Sinus mit drei Ziffern zu verwenden, die etwa von Grad zu Grad auf einem kleinen Blatt verzeichnet sind.

